\newpage
\chapter{Ответы на контрольные вопросы}

\begin{enumerate}
    \item \textbf{Что такое аналоговый сигнал? Какие параметры его характеризуют?} \\
    \textbf{Ответ:} Непрерывный во времени физический процесс, мгновенные значения которого могут принимать любые величины в непрерывном диапазоне.
    Характеризуется амплитудой $U$, частотой $f$, начальной фазой $\varphi_0$ и шириной спектра.

    \item \textbf{Какими операциями формируется цифровой сигнал? Объяснить каждую из них.} \\
    \textbf{Ответ:} Двумя операциями:
    \begin{itemize}
        \item \textbf{дискретизацией} --- замером значений непрерывного сигнала через равные дискретные интервалы времени $T_{\text{д}}$;
        \item \textbf{квантованием} --- округлением непрерывных значений отсчетов до ближайших разрешенных уровней с последующим кодированием в бинарный код.
    \end{itemize}

    \item \textbf{Почему низкочастотный информационный сигнал нельзя излучать передающей антенной напрямую в пространство?} \\
    \textbf{Ответ:} Из-за огромных требуемых габаритов антенных систем, ведь размер антенны соизмерим с длиной волны $\lambda = c/f$, а для частоты 1~кГц $\lambda = 300\text{~км}$ и невозможности одновременной передачи множества сигналов без взаимных помех.

    \item \textbf{С помощью какого преобразования формируется аналоговый сигнал из цифрового?} \\
    \textbf{Ответ:} С помощью цифро-аналогового преобразования --- ЦАП и сглаживающей фильтрации.

    \item \textbf{Дайте определение коэффициента амплитудной модуляции $m$. Каков его физический смысл?} \\
    \textbf{Ответ:} Коэффициент амплитудной модуляции --- отношение размаха изменения огибающей к удвоенной амплитуде несущей: $m = (U_{\max} - U_{\min})/(U_{\max} + U_{\min})$.
    Физический смысл --- показывает глубину модуляции несущего колебания информационным сигналом.

    \item \textbf{Что такое «перемодуляция» АМ-сигнала? К каким спектральным и временным последствиям для радиосигнала она приводит?} \\
    \textbf{Ответ:} Перемодуляция АМ-сигнала --- режим, при котором $m > 1$, другими словами, огибающая падает ниже нуля.
    Приводит к искажению формы сигнала и резкому неконтролируемому расширению спектра.

    \item \textbf{В чем заключается фундаментальное различие между частотной и фазовой модуляцией, если в обоих случаях управляемым параметром является полная фаза несущего колебания?} \\
    \textbf{Ответ:} При фазовой модуляции мгновенная фаза изменяется пропорционально самому первичному сигналу $s(t)$, а мгновенная частота пропорциональна его производной.
    При частотной модуляции мгновенная частота пропорциональна первичному сигналу, а полная фаза пропорциональна его интегралу $\int s(t)dt$.

    \item \textbf{В каких случаях полоса ЧМ существенно шире полосы АМ при одинаковой частоте модулирующего тона?} \\
    \textbf{Ответ:} При больших индексах модуляции ($m \gg 1$), когда девиация частоты $\Delta f_M$ значительно превышает частоту модулирующего тона $F$.

    \item \textbf{В чем заключается физическое различие между понятиями «символьная скорость» и «информационная скорость»? В каких единицах Международной системы (СИ) измеряется каждая из них?} \\
    \textbf{Ответ:} Символьная скорость — количество изменений сигнального состояния (символов) в секунду, измеряется в Бодах (в СИ --- $\text{с}^{-1}$).
    Информационная скорость — количество переданных бит данных в секунду; измеряется в \textbf{бит/с}.
    Связь: $R = R_s \cdot \log_2(M)$.

    \item \textbf{Что такое сигнальное созвездие на комплексной плоскости? Объясните назначение синфазной ($I$) и квадратурной ($Q$) осей при векторном анализе цифровых сигналов.} \\
    \textbf{Ответ:} Сигнальное созвездие --- графическое отображение дискретных состояний радиосигнала на комплексной плоскости.
    Ось $I$ --- синфазная, задает амплитуду косинусной составляющей несущей;
    ось $Q$ --- квадратурная, задает амплитуду синусной составляющей.

    \item \textbf{Почему при переходе от двухпозиционной фазовой модуляции ФМ-2 к четырехпозиционной ФМ-4 информационная скорость потока данных удваивается, а занимаемая полоса частот в эфире остается неизменной?} \\
    \textbf{Ответ:} В ФМ-4 один символ переносит 2 бита за счет 4 дискретных фазовых состояний ($M=4$).
    Символьная скорость $R_s$ и форма радиоимпульсов остаются неизменными, но информационный поток удваивается.

    \item \textbf{В чем заключается теоретическое преимущество квадратурной амплитудной модуляции n-QAM перед любой многопозиционной фазовой манипуляцией n-PSK при высокой плотности созвездия (от 16 точек и выше)?} \\
    \textbf{Ответ:} В n-PSK все точки расположены на одной окружности, и с ростом $n$ эвклидово расстояние между соседними точками катастрофически уменьшается, ухудшая помехоустойчивость.
    В n-QAM точки распределены по плоскости $I/Q$ равномерно сеткой, что дает максимальное расстояние между символами при заданной средней мощности.

    \item \textbf{Как изменится вид панорамы спектра цифрового сигнала на экране анализатора, если длительность одного символа уменьшить в 2 раза? Как это повлияет на тактовую частоту манипуляции и занимаемую полосу частот в эфире?} \\
    \textbf{Ответ:} Уменьшение длительности символа $T_s$ в 2 раза означает удвоение тактовой частоты манипуляции: $R_s = 1/T_s$.
    Главный лепесток спектра расширится в 2 раза, а занимаемая полоса частот в эфире пропорционально удвоится.
\end{enumerate}